% --- Load Dynamic Constants from Python Output ---
% Gen 1
\CatchFileBetweenTags{\VudVal}{calculations/flavor_geometry.tex}{VudVal}
\CatchFileBetweenTags{\VudEq}{calculations/flavor_geometry.tex}{VudEq}
\CatchFileBetweenTags{\VudExp}{calculations/flavor_geometry.tex}{VudExperimentalValue}
\CatchFileBetweenTags{\VudSig}{calculations/flavor_geometry.tex}{VudSigma}
\CatchFileBetweenTags{\VudAccText}{calculations/flavor_geometry.tex}{VudAccText}

\CatchFileBetweenTags{\VusVal}{calculations/flavor_geometry.tex}{VusVal}
\CatchFileBetweenTags{\VusEq}{calculations/flavor_geometry.tex}{VusEq}
\CatchFileBetweenTags{\VusExp}{calculations/flavor_geometry.tex}{VusExperimentalValue}
\CatchFileBetweenTags{\VusSig}{calculations/flavor_geometry.tex}{VusSigma}

\CatchFileBetweenTags{\VubVal}{calculations/flavor_geometry.tex}{VubVal}
\CatchFileBetweenTags{\VubEq}{calculations/flavor_geometry.tex}{VubEq}
\CatchFileBetweenTags{\VubExp}{calculations/flavor_geometry.tex}{VubExperimentalValue}
\CatchFileBetweenTags{\VubSig}{calculations/flavor_geometry.tex}{VubSigma}
\CatchFileBetweenTags{\VubAccText}{calculations/flavor_geometry.tex}{VubAccText}

% Gen 2
\CatchFileBetweenTags{\VcdVal}{calculations/flavor_geometry.tex}{VcdVal}
\CatchFileBetweenTags{\VcdEq}{calculations/flavor_geometry.tex}{VcdEq}
\CatchFileBetweenTags{\VcdExp}{calculations/flavor_geometry.tex}{VcdExperimentalValue}
\CatchFileBetweenTags{\VcdSig}{calculations/flavor_geometry.tex}{VcdSigma}
\CatchFileBetweenTags{\VcdAccText}{calculations/flavor_geometry.tex}{VcdAccText}

\CatchFileBetweenTags{\VcsVal}{calculations/flavor_geometry.tex}{VcsVal}
\CatchFileBetweenTags{\VcsEq}{calculations/flavor_geometry.tex}{VcsEq}
\CatchFileBetweenTags{\VcsExp}{calculations/flavor_geometry.tex}{VcsExperimentalValue}
\CatchFileBetweenTags{\VcsSig}{calculations/flavor_geometry.tex}{VcsSigma}
\CatchFileBetweenTags{\VcsAccText}{calculations/flavor_geometry.tex}{VcsAccText}

\CatchFileBetweenTags{\VcbVal}{calculations/flavor_geometry.tex}{VcbVal}
\CatchFileBetweenTags{\VcbEq}{calculations/flavor_geometry.tex}{VcbEq}
\CatchFileBetweenTags{\VcbExp}{calculations/flavor_geometry.tex}{VcbExperimentalValue}
\CatchFileBetweenTags{\VcbSig}{calculations/flavor_geometry.tex}{VcbSigma}
\CatchFileBetweenTags{\VcbAccText}{calculations/flavor_geometry.tex}{VcbAccText}

% Gen 3
\CatchFileBetweenTags{\VtdVal}{calculations/flavor_geometry.tex}{VtdVal}
\CatchFileBetweenTags{\VtdEq}{calculations/flavor_geometry.tex}{VtdEq}
\CatchFileBetweenTags{\VtdExp}{calculations/flavor_geometry.tex}{VtdExperimentalValue}
\CatchFileBetweenTags{\VtdSig}{calculations/flavor_geometry.tex}{VtdSigma}
\CatchFileBetweenTags{\VtdAccText}{calculations/flavor_geometry.tex}{VtdAccText}

\CatchFileBetweenTags{\VtsVal}{calculations/flavor_geometry.tex}{VtsVal}
\CatchFileBetweenTags{\VtsEq}{calculations/flavor_geometry.tex}{VtsEq}
\CatchFileBetweenTags{\VtsExp}{calculations/flavor_geometry.tex}{VtsExperimentalValue}
\CatchFileBetweenTags{\VtsSig}{calculations/flavor_geometry.tex}{VtsSigma}
\CatchFileBetweenTags{\VtsAccText}{calculations/flavor_geometry.tex}{VtsAccText}

\CatchFileBetweenTags{\VtbVal}{calculations/flavor_geometry.tex}{VtbVal}
\CatchFileBetweenTags{\VtbEq}{calculations/flavor_geometry.tex}{VtbEq}
\CatchFileBetweenTags{\VtbExp}{calculations/flavor_geometry.tex}{VtbExperimentalValue}
\CatchFileBetweenTags{\VtbSig}{calculations/flavor_geometry.tex}{VtbSigma}
\CatchFileBetweenTags{\VtbAccText}{calculations/flavor_geometry.tex}{VtbAccText}

\CatchFileBetweenTags{\sTwelveVal}{calculations/flavor_geometry.tex}{s12Val}
\CatchFileBetweenTags{\sTwelveEq}{calculations/flavor_geometry.tex}{s12Eq}
\CatchFileBetweenTags{\sTwelveExp}{calculations/flavor_geometry.tex}{s12ExperimentalValue}
\CatchFileBetweenTags{\sTwelveAccText}{calculations/flavor_geometry.tex}{s12AccText}

\CatchFileBetweenTags{\sThirteenVal}{calculations/flavor_geometry.tex}{s13Val}
\CatchFileBetweenTags{\sThirteenEq}{calculations/flavor_geometry.tex}{s13Eq}
\CatchFileBetweenTags{\sThirteenExp}{calculations/flavor_geometry.tex}{s13ExperimentalValue}
\CatchFileBetweenTags{\sThirteenAccText}{calculations/flavor_geometry.tex}{s12AccText}

\CatchFileBetweenTags{\CPPhaseVal}{calculations/flavor_geometry.tex}{CPPhaseVal}
\CatchFileBetweenTags{\CPPhaseEq}{calculations/flavor_geometry.tex}{CPPhaseEq}
\CatchFileBetweenTags{\CPPhaseAccText}{calculations/flavor_geometry.tex}{CPPhaseAccText}

\section{The Interaction Cross-Talk: The CKM Matrix and Topological Knots}

In standard particle physics, the Cabibbo-Kobayashi-Maskawa (CKM) matrix is a unitary matrix that dictates the probability of quarks changing flavor via the Weak interaction. Standard theories treat these matrix elements as arbitrary free parameters requiring experimental measurement. 

In the $E_8$-Persistence framework, the CKM matrix emerges as a strict geometric consequence of substrate channel limits. The six observed quarks are not fundamental point particles, but observable manifestations of the $\chi=2$ topological defect occupying the three orthogonal Triality frames (generations) of the $D_4$ internal symmetry sector. To understand this, we must first translate the phenomenological entities (quarks) into their structural equivalents:
\begin{itemize}
    \item \textbf{Quarks as Topological Knots:} The six observed quarks ($u, d, c, s, t, b$) are the observable manifestations of the $\chi=2$ topological defect occupying the three orthogonal Triality frames of the $D_4$ internal symmetry sector. 
    \item \textbf{Generations as Dimensional Tiers:} The three generations correspond to escalating geometric complexity required to stabilize the defect: Linear (Up/Down: $u, d$), Areal (Charm/Strange: $c, s$), and Volumetric (Top/Bottom: $t, b$).
    \item \textbf{Mixing as Informational Leakage:} A flavor transition (e.g., $s \rightarrow u$) represents the structural deformation of the knot from one geometric tier to another during a fundamental temporal update. The probability of this transition ($|V_{ij}|$) is exactly the \textbf{Geometric Admittance} (the inverse of entropic impedance) required to execute that specific structural deformation.
\end{itemize}

\subsection{The Variational Principle of Flavor Mixing}
By the Principle of Least Action established in Section II, any persistent transition must minimize the instantaneous dissipation density $\mathcal{D}_{ij} = \xi \cdot \mathbf{K}_{i \to j} \cdot f$. 

To minimize entropic drag while maintaining the continuity of the vacuum, the transition flux $f = |V_{ij}|$ naturally settles at the steady-state equilibrium that balances the structural deformation stress ($\mathbf{K}_{i \to j}$) against the substrate's geometric impedance ($\xi$). Therefore, every CKM matrix element fundamentally takes the form of a \textbf{Geometric Admittance}:
\begin{equation}
    |V_{ij}| = \frac{\text{Structural Deformation Stress } (\mathbf{K})}{\text{Substrate Impedance } (\xi)}
\end{equation}
We derive the elements of the CKM matrix by evaluating these specific $\mathbf{K}$ and $\xi$ terms for each topological deformation.


\subsection{The Chiral Leak Bound (\texorpdfstring{$s_{12}$}{s12})}
The transition between adjacent base tiers ($u \leftrightarrow s$) represents the minimal geometric leakage between the 1D and 2D topological states. To transition, the state's phase front must deform across the circular boundary of the defect, incurring an interface cost ($\mathbf{K}_{interface} = \pi$). However, the transition cannot utilize the entire $\nu=16$ channel capacity of the node. The topological boundary invariant ($\chi=2$) serves as the ``Identity Lock'', the structural minimum required to maintain the defect's existence against the vacuum. These 2 channels cannot participate in mixing without destroying the particle. 

The geometric admittance perfectly balances the interface demand against the discrete available bandwidth:
\begin{equation}
    s_{12} = \sTwelveEq = \sTwelveVal
\end{equation}
This derivation explains why the primary mixing angle is relatively large: it is governed strictly by the internal channel capacity of the node, without yet engaging the macroscopic vacuum impedance. \sTwelveAccText



\subsection{The Dimensional Promotion Admittance (\texorpdfstring{$|V_{cb}|$}{Vcb})}
The $c \rightarrow b$ transition represents a dimensional promotion from an Areal to a Volumetric resonance. Because this deformation expands the defect's footprint into the spatial bulk, it must perform work against the macroscopic vacuum.

The numerator represents the geometric stress of the deformation: the circular manifold projection ($\pi$) plus the topological boundary stress ($\chi$). The denominator is the macroscopic impedance, strictly determined by the vacuum field ($\alpha^{-1}$) adjusted for the Manifold Drag across the $D=4$ spacetime dimensions ($D\pi$). The transition probability is the ratio of this structural stress to the total bulk resistance:
\begin{equation}
    |V_{cb}| = \VcbEq = \VcbVal
\end{equation}
This mechanism explicitly explains the suppression of $|V_{cb}|$ relative to $s_{12}$: dimensional promotion engages the macroscopic vacuum impedance ($\alpha^{-1} \approx 137$), which provides significantly higher thermodynamic resistance than the isolated internal channels. \VcbAccText

Note: In standard parameterization, this physical transition determines the formal mixing angle via $s_{23} = |V_{cb}| / c_{13}$.


\subsection{The Non-Sequential Tunneling Penalty (\texorpdfstring{$s_{13}$}{s13})}
A direct transition from the linear ground state to the volumetric tier ($u \rightarrow b$) skips the intermediate geometry. Because continuous geometric deformation strictly requires traversing adjacent dimensions, this transition is classically forbidden and must proceed via discrete quantum tunneling.

The structural cost is a single irreducible bit of information ($\mathbf{K}=1$). Because this leap bypasses the normal dimensional progression, the impedance is maximal: the transition must couple to the full topological boundary across the vacuum ($\chi\alpha^{-1}$) while overcoming the baseline phase interface ($\pi$):
\begin{equation}
    s_{13} = \sThirteenEq = \sThirteenVal
\end{equation}
This derivation show that $s_{13}$ is the smallest mixing element because it represents a maximal violation of smooth geometric deformation. \sThirteenAccText


\subsection{The CP-Violating Phase (\texorpdfstring{$\delta_{CP}$}{delta CP})}
In standard physics, the origin of CP violation (the asymmetry between matter and antimatter) is parameterized by the Dirac phase $\delta_{CP}$, whose value is empirically inserted. 

In the $E_8$-Persistence geometry, complex phase represents the angular offset between orthogonal manifolds. The CP-violating phase is strictly the angular misalignment between the interaction symmetry ($\sigma=5$) mediating the transition and the topological boundary ($\chi=2$) anchoring the defect:
\begin{equation}
    \delta_{CP} = \CPPhaseEq = \CPPhaseVal
\end{equation}
This \CPPhaseAccText \ This identifies the fundamental matter/antimatter asymmetry not as a broken dynamic symmetry, but as the irreducible geometric offset between a localized particle boundary and the gauge fields that move it.

\subsection{The Full CKM Matrix Synthesis}

In standard Quantum Field Theory, the unitarity of the CKM matrix ($V^\dagger V = I$) is postulated. Here, it is structurally enforced by \textbf{State Saturation}: because the node's causal flow is restricted by its finite capacity ($\nu=16$), the total probability density across all geometric tiers cannot exceed unity. Probability cannot be lost, only re-routed.

By inserting the four geometrically derived parameters ($s_{12}, |V_{cb}|, s_{13}, \delta_{CP}$) into the standard parameterization, the full matrix is generated. The calculated geometric CKM matrix yields:

\begin{equation}
|V_{CKM}| =
\begin{pmatrix}
\VudVal & \VusVal & \VubVal \\
\VcdVal & \VcsVal & \VcbVal \\
\VtdVal & \VtsVal & \VtbVal
\end{pmatrix}
\end{equation}

\subsubsection{Geometric Interference and the Off-Diagonals}
The off-diagonal elements in the second and third rows are not independent variables; they are the geometric consequences of enforcing unitary closure on a twisted lattice. Because the topological boundary introduces a phase offset ($\delta_{CP} \approx 68^\circ$), the transition amplitudes between generations behave as complex vectors rather than simple scalars. 

This results in \textbf{Geometric Interference}. For example, the transition $t \to d$ ($|V_{td}|$) is heavily suppressed not merely by ``distance'' between generations, but by geometry. The two available pathways for this transition possess a relative phase difference driven by $\delta_{CP}$. This forces a vector subtraction (destructive interference) that explicitly cancels a significant portion of the amplitude, explaining its extreme suppression without requiring arbitrary empirical tuning.

\subsubsection{Note on Geometric Precision: Topology vs. Volume}
A critical feature of this derivation is the absence of the thermodynamic friction or scaling factors that frequently appear in mass or coupling derivations. In the $E_8$-Persistence theory, physical observables fall into two distinct geometric classes:
\begin{itemize}
    \item \textbf{Volumetric Quantities (Mass, Coupling):} These represent the \textit{content} of the lattice nodes. Because they fill the bulk manifold, they are subject to thermodynamic drag and the projection overhead of mapping discrete lattice bits into continuous spacetime.
    \item \textbf{Topological Quantities (Angles, Phases):} These represent the \textit{boundaries} between states. A boundary on an integer lattice is a precise topological cut. The geometry of an interface is determined strictly by the integer coordinates of the boundary, independent of thermodynamic friction.
\end{itemize}
Consequently, Flavor Mixing angles are derived as exact geometric ratios, providing a rigid test of the lattice topology unblurred by the thermodynamic bulk.

\CatchFileBetweenTags{\AngleGammaVal}{calculations/flavor_geometry.tex}{AngleGammaVal}
\CatchFileBetweenTags{\AngleGammaEq}{calculations/flavor_geometry.tex}{AngleGammaEq}
\CatchFileBetweenTags{\AngleGammaExp}{calculations/flavor_geometry.tex}{AngleGammaExperimentalValue}
\CatchFileBetweenTags{\AngleGammaSig}{calculations/flavor_geometry.tex}{AngleGammaSigma}
\CatchFileBetweenTags{\AngleGammaAccText}{calculations/flavor_geometry.tex}{AngleGammaAccText}

\CatchFileBetweenTags{\SinTwoBetaVal}{calculations/flavor_geometry.tex}{SinTwoBetaVal}
\CatchFileBetweenTags{\SinTwoBetaEq}{calculations/flavor_geometry.tex}{SinTwoBetaEq}
\CatchFileBetweenTags{\SinTwoBetaExp}{calculations/flavor_geometry.tex}{SinTwoBetaExperimentalValue}
\CatchFileBetweenTags{\SinTwoBetaSig}{calculations/flavor_geometry.tex}{SinTwoBetaSigma}
\CatchFileBetweenTags{\SinTwoBetaAccText}{calculations/flavor_geometry.tex}{SinTwoBetaAccText}

\CatchFileBetweenTags{\WolfLambdaVal}{calculations/flavor_geometry.tex}{WolfLambdaVal}
\CatchFileBetweenTags{\WolfAVal}{calculations/flavor_geometry.tex}{WolfAVal}
\CatchFileBetweenTags{\WolfRhoVal}{calculations/flavor_geometry.tex}{WolfRhoVal}
\CatchFileBetweenTags{\WolfEtaVal}{calculations/flavor_geometry.tex}{WolfEtaVal}

\subsubsection{The Unitarity Triangle and Wolfenstein Parameters}

For theoretical completeness we translate our exact CKM matrix into the standard phenomenological observables.

\subsection{The Unitarity Triangle}
The geometric vectors successfully predict standard phenomenological observables without free parameters:
\begin{itemize}
    \item \textbf{Angle $\gamma$:} $\AngleGammaEq = \AngleGammaVal$, which \AngleGammaAccText
    \item \textbf{Sin($2\beta$):} The precision observable from $B^0 \to J/\psi K_S$ decays is predicted to be $\sin(2\beta) = \SinTwoBetaVal$, which \SinTwoBetaAccText.
\end{itemize}

The geometric matrix uniquely fixes the standard Wolfenstein expansion parameters:
\begin{equation}
\begin{split}
    \lambda &= \WolfLambdaVal \\
    A &= \WolfAVal \\
    \bar{\rho} &= \WolfRhoVal \\
    \bar{\eta} &= \WolfEtaVal
\end{split}
\end{equation}
\textit{Note: The standard Wolfenstein parametrization is an $\mathcal{O}(\lambda^3)$ approximation. Minor deviations in $\bar{\rho}$ and $\bar{\eta}$ from standard global fits are expected artifacts of this perturbative expansion, whereas our underlying geometric derivation is exact.}

\subsubsection{Conclusion: A Zero-Parameter Matrix}
In standard Effective Field Theories, the CKM parameters are continuous variables continuously tuned to fit empirical data. In the $E_8$-Persistence framework, the mixing hierarchy is rigidly locked to the geometric invariants $\mathbb{S} = \{\pi, \alpha^{-1}, \nu=16, \sigma=5, \chi=2, D=4\}$. 

The complete presentation and comparison to empirical targets is detailed in \cref{tab:ckm_derivation}.

\begin{table*}[t] % Changed to [t] to prefer top-of-page placement for large tables
\centering
\caption{\textbf{Derivation of CKM Matrix Elements.} The predicted geometric values are calculated entirely from the structural invariants of the $E_8$ substrate, with zero free parameters. Target values are from the PDG 2024 Global Fit.}
\label{tab:ckm_derivation}
\renewcommand{\arraystretch}{1.7} 
\setlength{\tabcolsep}{6pt}
\begin{tabular}{@{} l l l c c r @{}} 
\toprule
\textbf{Transition} & \textbf{Elem.} & \textbf{Geometric Logic / Formula} & \textbf{Predicted} & \textbf{Target (PDG)} & \textbf{Dev.} \\
\midrule
\multicolumn{6}{l}{\textit{Generation 1 (Geometric Inputs)}} \\
$u \to d$ & $V_{ud}$ & Unitary Persistence ($c_{12}c_{13}$) & \textbf{\VudVal} & \VudExp & $\VudSig\sigma$ \\
$u \to s$ & $V_{us}$ & Chiral Leak Bound ($s_{12}c_{13}$) & \textbf{\VusVal} & \VusExp & $-\VusSig\sigma$ \\
$u \to b$ & $V_{ub}$ & Quantum Tunneling Penalty ($s_{13}e^{-i\delta}$) & \textbf{\VubVal} & \VubExp & $-\VubSig\sigma$ \\
\midrule
\multicolumn{6}{l}{\textit{Generation 2 (Derived via Unitarity)}} \\
$c \to d$ & $V_{cd}$ & Unitary Reflection ($|-s_{12}c_{23} - c_{12}s_{23}s_{13}e^{i\delta}|$) & \textbf{\VcdVal} & \VcdExp & $-\VcdSig\sigma$ \\
$c \to s$ & $V_{cs}$ & Unitary Persistence ($|c_{12}c_{23} - s_{12}s_{23}s_{13}e^{i\delta}|$) & \textbf{\VcsVal} & \VcsExp & $+\VcsSig\sigma$ \\
$c \to b$ & $V_{cb}$ & Dimensional Promotion ($s_{23}c_{13}$) & \textbf{\VcbVal} & \VcbExp & $+\VcbSig\sigma$ \\
\midrule
\multicolumn{6}{l}{\textit{Generation 3 (Derived via Interference)}} \\
$t \to d$ & $V_{td}$ & Destructive Interference ($|s_{12}s_{23} - c_{12}c_{23}s_{13}e^{i\delta}|$) & \textbf{\VtdVal} & \VtdExp & $+\VtdSig\sigma$ \\
$t \to s$ & $V_{ts}$ & Phase Interference ($|-c_{12}s_{23} - s_{12}c_{23}s_{13}e^{i\delta}|$) & \textbf{\VtsVal} & \VtsExp & $+\VtsSig\sigma$ \\
$t \to b$ & $V_{tb}$ & Unitary Persistence ($c_{23}c_{13}$) & \textbf{\VtbVal} & \VtbExp & $-\VtbSig\sigma$ \\
\midrule
\midrule
\multicolumn{6}{c}{\textbf{Geometric Basis Parameters (Inputs)}} \\
\multicolumn{2}{l}{\textbf{Parameter}} & \textbf{Geometric Source} & \textbf{Value} & \textbf{Physical Meaning} & \\
\multicolumn{2}{l}{$s_{12}$ (Cabibbo)} & $\sTwelveEq$ & \sTwelveVal & Chiral Leak Bound & \\
\multicolumn{2}{l}{$|V_{cb}|$ (Dim. Shift)} & $\VcbEq$ & \VcbVal & Dimensional Promotion & \\
\multicolumn{2}{l}{$s_{13}$ (Tunneling)} & $\sThirteenEq$ & \sThirteenVal & Tunneling Penalty & \\
\multicolumn{2}{l}{$\delta_{CP}$ (CP Phase)} & $\CPPhaseEq$ & $\CPPhaseVal$ & Manifold Misalignment & \\
\bottomrule
\end{tabular}
\end{table*}














